\documentclass[10pt]{article}

\usepackage[margin=0.8in]{geometry}
\usepackage{booktabs}
\usepackage{hyperref}
\usepackage{enumitem}
\usepackage{amsmath}
\usepackage{microtype}
\usepackage{xcolor}
\usepackage{array}

\hypersetup{
  colorlinks=true,
  linkcolor=blue,
  citecolor=blue,
  urlcolor=blue
}

\title{Bio-Inspired Inference Control for Local Large Language Models:\\
A DNA Gene-Chain Architecture in Rust}

\author{
  Hao Yang YangHao\\
  Independent Researcher\\
  \texttt{2499510083@qq.com}\\
  \url{https://claude.chiddns.com}\\
  Code: \url{https://github.com/yanghao1143/rust-norion}\\
  Draft technical report, June 2026
}

\date{}

\begin{document}
\maketitle

\begin{abstract}
Long-context large language model (LLM) applications increasingly require more
than a fast forward kernel. Private deployments, multi-turn coding assistants,
tool-using agents, and local research systems must decide which context to keep,
which memory to persist, when to revise a control policy, and how to reject
unsafe self-modification. This report presents the architecture of
\texttt{rust-norion}, a Rust prototype for a bio-inspired inference
control-layer engine. The system treats reusable reasoning context as a
``reasoning genome'' composed of explicit expression segments, latent memory
segments, mutation evidence, rollback anchors, and writer-gated evolution
records. Unlike runtime-first engines that primarily optimize model execution,
\texttt{rust-norion} focuses on the control plane around inference:
disk-backed KV memory, hierarchical routing, adaptive attention thresholds,
reinforced KV-Fusion, task-aware hierarchy adjustment, reflection loops,
gene-chain splicing, mutation isolation, and auditable self-goal planning. The
current implementation is a prototype: durable genome, memory, experiment, and
self-goal writes remain preview-only unless explicit gates, validation evidence,
rollback plans, and operator approval are present. We describe the motivation,
system model, core abstractions, governance gates, current implementation
status, and evaluation plan needed to turn bio-inspired control from metaphor
into measurable engineering behavior.
\end{abstract}

\noindent\textbf{Keywords:} large language models, inference control, local AI,
long context, KV cache, self-evolution, bio-inspired architecture, Rust.

\section*{Artifact Availability}

The project code is publicly available at
\url{https://github.com/yanghao1143/rust-norion}. This draft describes the
\texttt{codex-runtime-device-abi} development branch as of June 2026. For an
external arXiv submission, the authors should create a stable release tag or
archived commit snapshot and ensure that the repository license statement
matches the intended reuse policy at submission time. The current draft is
written as a compact 6--8 page systems-style technical report; exact page count
depends on the final TeX engine and reference formatting.

\section{Introduction}

The dominant open-source LLM infrastructure stack separates into two broad
families. Runtime-first systems focus on model loading, kernels, quantization,
streaming, batching, and hardware utilization. Gateway-first systems focus on
request routing, rate limits, observability, and policy. Both layers are
necessary, but neither fully addresses a third problem: how a local model
should govern its own context, memory, routing thresholds, reflection outputs,
and bounded self-improvement over long periods without frequent weight
retraining.

\texttt{rust-norion} explores this third layer. It is not designed as a direct
competitor to high-throughput inference runtimes. Instead, it is a Rust
control-plane prototype that can sit around a deterministic local runtime today
and around stronger pluggable model backends later. The design assumes that
many practical improvements in local AI behavior come from better inference
control rather than from immediately changing model weights: selecting useful
context, pruning redundant segments, reusing KV memory, adjusting hierarchy
weights by task profile, learning from reflection, and blocking unsafe durable
state changes.

The architecture is intentionally bio-inspired. A biological genome is not only
a flat sequence of expressed material; it contains regulatory structure,
non-coding regions, mutation handling, inheritance, repair, and irreversible
failure boundaries. \texttt{rust-norion} maps these ideas into software
primitives: an express chain for active reasoning context, a memory chain for
latent reusable state, a splicer for filtering low-value context, mutation
detectors for malformed or harmful segments, gene scissors for reversible edit
plans, and writer gates for previewing durable changes before application.

This report makes three contributions:

\begin{enumerate}[leftmargin=*]
  \item It defines a DNA-inspired inference control architecture that separates
  model runtime execution from memory, routing, context governance, and
  self-evolution policy.
  \item It describes a Rust implementation strategy for deterministic previews,
  digest-only evidence, rollback anchors, and explicit apply gates across
  memory, genome, experiment-ledger, and self-goal queues.
  \item It identifies a focused evaluation plan for measuring whether the
  architecture improves long-context efficiency, KV reuse, mutation isolation,
  and safe self-improvement without claiming production runtime performance.
\end{enumerate}

\section{Scope and Terminology}

Several terms in \texttt{rust-norion} are project-specific and should be read
as control-plane mechanisms rather than biological simulation claims.
\textbf{Noiron} denotes the local inference control loop around a model
runtime. \textbf{FHT-DKE} denotes the project's hierarchical routing and
disk-backed knowledge execution target: adaptive attention routing, sparse
context selection, KV exchange, and tiered persistent memory are planned and
validated through explicit Rust structures before they are delegated to a
runtime. \textbf{Reasoning genome} denotes persistent strategy state for
retrieval, routing, reflection, safety, and tool use. \textbf{Gene scissors}
denotes reversible edit plans for that strategy state. These mechanisms do not
change neural weights directly.

The architecture therefore sits between a model-serving runtime and an
application. It can be evaluated even before a production model kernel exists,
because many failure modes are control-layer failures: context pollution,
private-data retention, stale memory, unstable thresholds, invalid KV shape,
unsafe tool policy, and unreviewed self-modification. This is why the current
Rust implementation emphasizes deterministic previews, trace schemas, writer
gates, rollback anchors, and focused benchmarks.

\section{Background and Motivation}

Long-context inference exposes a resource-management problem. StreamingLLM
showed that retaining a small set of attention-sink tokens can stabilize
streaming language modeling beyond the nominal context window
\cite{streamingllm}. CEPE shows that parallel context encoding can extend
decoder-only models by chunking long inputs and exposing additional context via
cross-attention \cite{cepe}. Biosequence models such as Evo demonstrate that
Transformer-style sequence modeling can represent biological structure across
molecular and genome scales \cite{evo}. These works motivate, but do not by
themselves provide, a local software control plane for deciding how an LLM
should persist, splice, repair, inherit, and audit its own inference context.

The \texttt{rust-norion} hypothesis is that long-running local agents need an
explicit inference genome: a structured control state that records what should
be expressed now, what should remain latent, what can be reused, what should be
cut, and which changes are too risky to apply automatically. This is especially
important for local coding assistants and enterprise-style deployments where
raw prompts, private datasets, credentials, model weights, and unreviewed
third-party code must not leak into persistent memory.

The project therefore optimizes for a different target than a runtime benchmark
leaderboard. The desired property is a safe adaptive loop:

\[
  \text{observe} \rightarrow \text{route} \rightarrow \text{infer}
  \rightarrow \text{reflect} \rightarrow \text{propose}
  \rightarrow \text{validate} \rightarrow \text{apply or rollback}.
\]

At the current prototype stage, the last step is intentionally blocked for many
state classes. Preview reports can compute candidate writes and apply plans,
but durable side effects require explicit gates.

\section{System Overview}

\subsection{Control Plane and Runtime Boundary}

The system separates a model runtime from an inference control plane. The model
runtime owns tokenizer behavior, embeddings, architecture shape, weights,
forward kernels, and optional KV import/export. The control plane owns context
selection, memory tiering, routing thresholds, reflection, process reward,
recursive scheduling, genome edits, and durable-state governance.

This separation is implemented through stable Rust traits such as
\texttt{InferenceBackend} and \texttt{ModelRuntime}. A deterministic local
runtime can be used for tests and ABI validation. A production backend can later
attach behind the same boundary if it satisfies manifest, device-contract, KV
shape, token uncertainty, and trace requirements. The point is not to make the
prototype kernel outperform mature runtimes; the point is to keep the
control-layer logic portable across runtimes.

\begin{figure}[t]
\centering
\fbox{\begin{minipage}{0.94\linewidth}
\small
\textbf{Input and memory evidence}
$\rightarrow$ splicing and mutation detection
$\rightarrow$ express-chain projection
$\rightarrow$ adaptive routing and hierarchy planning
$\rightarrow$ model runtime request
$\rightarrow$ reflection and process reward
$\rightarrow$ preview-only memory/genome/self-goal proposal
$\rightarrow$ writer gate
$\rightarrow$ explicit apply or rollback.
\end{minipage}}
\caption{High-level closed loop. The runtime performs model execution; the
control plane governs what context is expressed, what memory is fused, what
state is persisted, and which proposed changes are blocked or admitted.}
\label{fig:closed-loop}
\end{figure}

\subsection{Reasoning Genome}

The reasoning genome is a structured representation of inference control state.
The main conceptual components are:

\begin{itemize}[leftmargin=*]
  \item \textbf{Express chain}: active segments that may affect the current
  response or route decision.
  \item \textbf{Memory chain}: latent reusable segments that may be retrieved,
  fused, decayed, or forgotten.
  \item \textbf{Gene segment}: a typed unit of context, evidence, routing
  policy, memory admission, or repair metadata.
  \item \textbf{Splicer}: a filter that separates useful expressed segments
  from intron-like redundant material.
  \item \textbf{Mutation detector}: a classifier for malformed, stale,
  conflicting, private, executable, or low-fitness segments.
  \item \textbf{Gene scissors}: a reversible edit planner that records before
  and after digests, forensic copies, stable anchors, lineage, and rollback
  conditions.
\end{itemize}

The biological analogy is constrained by engineering value. ``Intron'' means
low-value context that consumes budget. ``Mutation'' means a segment that is
format-invalid, harmful, stale, private, or behaviorally regressive. ``Repair''
means a validated replacement or isolation plan, not uncontrolled model
rewriting.

\begin{table}[t]
\centering
\small
\begin{tabular}{p{0.24\linewidth}p{0.31\linewidth}p{0.32\linewidth}}
\toprule
Biological inspiration & Software mapping & Engineering purpose \\
\midrule
Expressed sequence & \texttt{express\_chain} & Small active strategy
surface for routing, memory retrieval, reflection, and tool posture. \\
Non-coding or latent sequence & \texttt{memory\_chain} & Provenance,
validation, source digests, rollback anchors, and cold reusable evidence. \\
Splicing & Splicer and intron filter & Remove low-value context before it
enters expensive attention or KV prefill. \\
Variant or mutation & Mutation detector & Isolate malformed, stale, private,
conflicting, or low-fitness segments. \\
Repair and regeneration & Gene scissors plan & Relabel, quarantine, cut,
splice, regenerate, or roll back only through preview and validation gates. \\
Inheritance & Genome ledger and session state & Carry validated strategies
across runs while preserving lineage and rollback ability. \\
\bottomrule
\end{tabular}
\caption{DNA-inspired concepts are used as engineering metaphors with explicit
control-plane behavior.}
\label{tab:mapping}
\end{table}

\subsection{Disk-Backed KV and Memory Tiers}

\texttt{rust-norion} uses disk-backed memory plans and ledger-style write gates
so useful context can persist across sessions without turning raw prompts into
uncontrolled durable data. Hot/Warm/Cold tiering, KV-Fusion, and memory
consolidation are treated as control-plane decisions. Memory candidates carry
quality signals, process reward, privacy classification, rollback anchors,
review packets, and read-only/write/applied flags.

The design goal is to make local long-context behavior practical on constrained
hardware. Rather than caching every token forever, the control layer can prefer
high-fitness memory, demote cold state, and keep write admission explicit.

\subsection{Trace Surface and Evidence Hygiene}

Every adaptive system needs a trace surface that is useful for debugging but
safe for persistence. \texttt{rust-norion} treats traces as evidence packets,
not as raw transcript archives. Trace records should therefore contain schema
versions, stable ids, counters, decisions, redaction digests, budget summaries,
rollback anchors, and reason codes. They should not contain secrets, raw
private prompts, hidden chain-of-thought text, raw model weights, private
datasets, copied third-party source, or unreviewed executable payloads.

This choice is central to the self-evolution story. A system cannot safely
learn from prior runs if the learning substrate becomes a privacy sink. The
project's trace schema gates are designed to reject suspicious durable evidence
before it is used for memory admission, genome edits, experiment promotion, or
self-goal queue planning.

\section{Adaptive Routing and Hierarchy}

The routing layer estimates when a token or segment should use local attention,
global attention, projection, convolution-style fusion, memory lookup, or
fallback paths. Inputs include entropy, task profile, context length, hardware
pressure, KV budget, cache hit rate, latency pressure, and device profile.

Task-aware hierarchy adjustment stores profile-specific control weights. Coding
tasks may prioritize compiler/test evidence and precise local context. Long
document tasks may prioritize recursive summaries and stable anchors. General
reasoning may prioritize uncertainty and reflection feedback. In each case, the
control plane is allowed to adjust thresholds and hierarchy weights, but not
silently mutate model weights.

Reflection provides another feedback channel. Draft quality, structured issues,
revision actions, and process reward can update future routing and memory
admission. The target is test-time adaptation: improve how the model is used
without assuming that every interaction triggers training.

\subsection{Task-Scoped Adaptation}

The control layer keeps task profiles separate because a useful strategy for
one mode can be harmful in another. Rust coding tasks should elevate compiler,
test, benchmark, and clean-room evidence. Chinese-language interaction can
prefer bilingual summaries, language-mode continuity, and culturally stable
response style. Long-document tasks can prefer recursive scheduling,
hierarchical gist records, and segment-local repair. General reasoning can
prefer uncertainty, reflection diagnostics, and contradiction pressure. The
profile boundary prevents one successful local heuristic from overwriting the
global behavior of the engine.

This task-scoped design also gives the biological metaphor a measurable
interpretation. A ``phenotype'' is the observable behavior of a profile under a
benchmark or trace replay: latency, context budget, memory hit rate, reflection
issues, compiler success, or repair outcome. A ``fitness'' score is not a vague
preference; it is a weighted control-plane metric that can be checked before a
gene is promoted.

\section{Self-Evolution Governance}

\subsection{Preview-Only State Changes}

The project treats durable state mutation as a safety boundary. Memory writes,
genome edits, experiment-ledger promotions, and self-goal queue appends default
to preview-only. Preview reports can include counts, digests, rollback anchors,
and reason codes, but raw prompts, hidden reasoning, credentials, private
payloads, and unreviewed source content must not enter durable traces.

\subsection{Unified Writer Gate}

The unified writer gate normalizes candidate writes from multiple domains. It
evaluates review packet references, evidence references, rollback readiness,
privacy checks, license checks, operator approval, approval-reference matching,
source read-only flags, source write flags, and source applied flags. Its
decisions are:

\begin{itemize}[leftmargin=*]
  \item \texttt{preview\_only}: safe evidence exists, but durable writes are
  disabled.
  \item \texttt{hold}: required evidence is incomplete.
  \item \texttt{reject}: privacy, license, source-write, source-active, or
  source-applied safety failed.
  \item \texttt{ready\_for\_explicit\_apply}: all gates passed under an explicit
  write-enabled policy, but application still requires a scoped apply workflow.
\end{itemize}

\subsection{Self-Goal Queue}

The self-goal mechanism answers when the system can set its own goals. The
answer is staged. First, \texttt{SelfGoalProposalReport} proposes bounded
candidate goals. Second, \texttt{SelfGoalAdmissionGate} evaluates whether a
candidate would be admissible under success evidence, budget, rollback, and
approval constraints. Third, \texttt{SelfGoalQueuePreviewGate} emits a
digest-only append packet for one admissible candidate. Fourth, the unified
writer gate preflights that packet as an \texttt{evolution\_goal\_queue}
candidate. Fifth, \texttt{SelfGoalQueueApplyPlanner} builds a digest-only apply
plan with live queue digest, rollback anchor, append digest, expected resulting
queue digest, and matching writer-gate refs. Even here,
\texttt{write\_allowed=false} and \texttt{applied=false}; a future append
executor must apply the plan explicitly.

This design lets the system reason about its next work without giving it a
blank check to mutate durable project state.

\subsection{Why Autonomous Goals Remain Gated}

Autonomous goal generation is useful only if it stays bounded. In this
architecture, a self-goal is a candidate control-plane improvement, not an
instruction to rewrite arbitrary files. The candidate must pass proposal,
admission, queue preview, writer-gate preflight, and apply-plan checks. Each
stage reduces the risk of runaway behavior: stale queue digests are rejected,
duplicate goals are held, incomplete validation evidence blocks admission,
source records marked as already applied are rejected, and explicit operator
approval is required before durable queue mutation. This lets the engine plan
future evolution while preserving human ownership of irreversible changes.

\section{Implementation Status}

The public code is available at
\url{https://github.com/yanghao1143/rust-norion}. As of this draft, the
repository contains a Rust prototype with deterministic control-plane modules,
trace schema gates, local runtime boundaries, memory admission previews,
reasoning genome governance, gene-scissors transaction journals, self-evolution
preflights, writer gates, and self-goal queue apply planning. The implementation
also includes tests for preview-only behavior, digest-only records, rollback
anchors, unsafe source-flag rejection, duplicate-goal holds, and writer-gate
readiness.

The project should be interpreted as a control-layer prototype, not as a
production LLM runtime. It does not claim to outperform mature Rust inference
runtimes on kernel throughput. The current deterministic runtime is valuable
because it makes governance behavior testable and reproducible. Production
kernel integration remains a runtime-boundary problem.

The implementation is deliberately shaped as small Rust modules rather than a
single monolithic agent. Examples include \texttt{router.rs} for adaptive route
planning, \texttt{hierarchy.rs} for task-aware hierarchy weights,
\texttt{memory\_admission.rs} and \texttt{kv\_cache.rs} for KV governance,
\texttt{reasoning\_genome.rs} for genome records and gene-scissors transactions,
\texttt{writer\_gate.rs} for shared durable-write evaluation,
\texttt{self\_goal\_proposal.rs} for bounded self-goal planning, and
\texttt{trace.rs} for schema-gated evidence. These modules are exported through
\texttt{src/lib.rs} so examples, benches, and future services can exercise the
control plane without depending on private internals.

\begin{table}[t]
\centering
\small
\begin{tabular}{lll}
\toprule
Layer & Purpose & Current stance \\
\midrule
Runtime boundary & Backend traits, manifest, device contract & Prototype with deterministic paths \\
KV memory & Disk-backed admission and tier plans & Preview/gated durable writes \\
Reasoning genome & Gene segments, scissors, lineage, repair & Preview/gated edits \\
Routing/hierarchy & Adaptive thresholds and task profiles & Control-plane tests and traces \\
Self-evolution & Experiments, promotion, rollback & Preview/gated promotion \\
Self-goals & Proposal, admission, preview, apply plan & No autonomous append executor yet \\
\bottomrule
\end{tabular}
\caption{Summary of the rust-norion prototype layers.}
\end{table}

\section{Reproducibility and arXiv Readiness}

An arXiv-facing version of this report should make the artifact easy to inspect
and difficult to misunderstand. The recommended release process is:

\begin{enumerate}[leftmargin=*]
  \item Create a stable Git tag for the exact code snapshot discussed in the
  report and include that tag in the artifact section.
  \item Keep the TeX source self-contained. arXiv's TeX guidance recommends
  preparing the source package carefully and checking the generated PDF before
  final submission \cite{arxivtex}.
  \item Include a short command list for local verification, such as
  \texttt{cargo fmt --all -- --check}, focused \texttt{cargo test} invocations
  for routing, genome, writer-gate, and self-goal modules, and
  \texttt{cargo check}.
  \item Publish only benchmark numbers that have a reproducible script,
  baseline, after-change result, and rollback/failure interpretation.
  \item State the license boundary clearly. If the project moves to a
  permissive open core, the paper and repository must be updated together so
  readers do not confuse future policy with the current controlling license.
\end{enumerate}

The likely arXiv subject fit is computer science infrastructure or language
model systems, with \texttt{cs.CL} or \texttt{cs.AI} as possible categories
depending on the final emphasis. A pure position paper is weaker than a
systems-style technical report with code, traces, and small reproducible
experiments. The first public version should therefore prioritize a credible
architecture claim and executable evidence over broad performance claims.

\section{Evaluation Plan}

The architecture should be evaluated with measurements aligned to its real
claim: better inference control, not raw kernel speed. The proposed benchmark
suite should include:

\begin{enumerate}[leftmargin=*]
  \item \textbf{Context compression}: token or segment reduction from splicing
  while preserving task-relevant anchors.
  \item \textbf{Mutation isolation}: rate at which malformed/private/executable
  segments are isolated before expression or durable write.
  \item \textbf{Repair success}: proportion of structured segment errors that
  can be corrected and validated without contaminating the full chain.
  \item \textbf{KV reuse}: hit rate and latency improvement from reinforced
  KV-Fusion and tiered residency.
  \item \textbf{Routing efficiency}: reduction in wasted global attention or
  long-context compute under equivalent task quality.
  \item \textbf{Self-evolution safety}: count of proposed changes, admitted
  previews, held plans, rejected unsafe writes, applied changes, and rollbacks.
  \item \textbf{Local deployment reproducibility}: deterministic behavior across
  CPU-only, integrated GPU, mobile, edge, browser-WASM, and accelerator profile
  simulations where possible.
\end{enumerate}

The most important reporting rule is to keep baseline and after-change evidence
together. A self-evolution claim should not be accepted if it only shows that a
loop ran; it should show measurable control-plane movement, validation success,
and rollback safety.

\subsection{Minimum Demonstration Suite}

The first public demonstration can be intentionally small:

\begin{itemize}[leftmargin=*]
  \item a splicing demo that marks exon, intron, and variant segments and shows
  a context-compression ratio;
  \item a mutation-repair demo that detects stale labels, invalid metadata,
  missing source hashes, private markers, or KV-shape errors and emits a
  read-only repair plan;
  \item a memory admission demo that shows why a useful segment is accepted,
  held, or rejected before it can enter disk-backed KV;
  \item a self-goal demo that proposes one next improvement, admits it only
  under complete evidence, and stops at a digest-only apply plan;
  \item a rollback demo that proves an admitted control-plane change has a
  stable previous state.
\end{itemize}

This suite is modest, but it is enough to make the core claim testable: the
system can govern long-context memory and self-evolution as structured,
auditable control state.

\section{Open Source Strategy}

The project is best positioned as a small, distinctive research/engineering
codebase rather than a broad model-serving platform. Community work should
therefore focus on gene-chain control, memory governance, runtime boundaries,
tests, examples, and reproducible benchmark evidence.

For broader adoption, the long-term license direction should favor a permissive
open core where contributors and users can build commercial products around the
basic control-layer primitives, while enterprise-specific multi-tenancy,
advanced audit suites, private deployment support, and specialized evolution
operators can remain commercial add-ons. Until a formal license transition is
approved, the repository's existing license remains controlling.

\section{Threats to Validity}

The main technical risk is that control-plane metrics become proxies for actual
task quality. A lower context budget is not valuable if answer quality drops. A
higher cache hit rate is not valuable if stale memory is reused. A repair plan
is not valuable if it silently hides important evidence. Future experiments
must therefore pair efficiency measurements with correctness, safety, and
human-review signals.

The second risk is overfitting to the deterministic prototype runtime. The
control layer should keep deterministic tests because they are excellent for
regression gates, but the same policies must later be tested against stronger
runtime backends with realistic tokenizer behavior, KV layout, latency, and
failure modes.

The third risk is terminology drift. Biological language is useful for naming
structured control problems, but it can mislead readers if it implies biological
fidelity. The report therefore treats DNA, splicing, mutation, repair, and
rejuvenation as disciplined software analogies. The project succeeds only if
those analogies produce measurable control behavior.

\section{Limitations}

The architecture is early and should not be oversold. Several risks remain.
First, the biological metaphor can become harmful if it adds terminology
without measurable benefit. Every gene-chain primitive must correspond to a
clear engineering behavior. Second, a single-maintainer project cannot
reasonably compete with runtime teams on every model architecture, accelerator,
and optimized kernel. The backend boundary must stay pluggable. Third,
self-evolution is only useful if it is validated and reversible; otherwise it
becomes technical debt with a more exciting name. Fourth, privacy and licensing
governance must remain strict because long-lived memory can amplify mistakes.

\section{Conclusion}

\texttt{rust-norion} proposes a bio-inspired inference control architecture for
local LLM systems. The core idea is to treat context, memory, routing, repair,
and self-evolution as a governed reasoning genome around a pluggable model
runtime. This shifts the research question from ``How fast is the kernel?'' to
``How safely and efficiently can a local model use, persist, repair, and evolve
its control state?'' The current Rust prototype implements many of the
governance surfaces needed to explore that question: digest-only previews,
rollback anchors, gene-scissors journals, memory admission plans, writer gates,
and self-goal apply planning. The next step is to connect these surfaces to
focused demos and benchmarks that make the benefits measurable.

\begin{thebibliography}{9}

\bibitem{attention}
Ashish Vaswani, Noam Shazeer, Niki Parmar, Jakob Uszkoreit, Llion Jones,
Aidan N. Gomez, Lukasz Kaiser, and Illia Polosukhin. Attention Is All You Need.
\textit{Advances in Neural Information Processing Systems}, 2017.

\bibitem{streamingllm}
Guangxuan Xiao, Yuandong Tian, Beidi Chen, Song Han, and Mike Lewis. Efficient
Streaming Language Models with Attention Sinks. arXiv:2309.17453, 2023.
\url{https://arxiv.org/abs/2309.17453}.

\bibitem{cepe}
Howard Yen, Tianyu Gao, and Danqi Chen. Long-Context Language Modeling with
Parallel Context Encoding. arXiv:2402.16617, 2024.
\url{https://arxiv.org/abs/2402.16617}.

\bibitem{evo}
Eric Nguyen et al. Sequence modeling and design from molecular to genome scale
with Evo. \textit{Science}, 2024.
\url{https://www.science.org/doi/abs/10.1126/science.ado9336}.

\bibitem{evo-code}
Arc Institute and collaborators. Evo: biological foundation modeling from
molecular to genome scale. \url{https://github.com/evo-design/evo}, 2024--2026.

\bibitem{arxivtex}
arXiv. Submit TeX/LaTeX. arXiv help documentation, accessed June 2026.
\url{https://info.arxiv.org/help/submit_tex.html}.

\bibitem{rustnorion}
yanghao1143 and contributors. rust-norion: DNA-inspired inference control layer
in Rust. \url{https://github.com/yanghao1143/rust-norion}, 2026.

\end{thebibliography}

\end{document}
